\documentclass[12pt]{article}
\usepackage{amsmath, amsfonts, amssymb, amsthm, fullpage, color, bbm, enumerate}
\usepackage[T1]{fontenc}
\usepackage{lmodern}

\usepackage[onehalfspacing]{setspace}

\usepackage[hidelinks]{hyperref}
%\usepackage[capitalize,nameinlink,noabbrev]{cleveref}

\theoremstyle{plain}
\newtheorem{asn}{Assumption}
%\crefname{asn}{Assumption}{Assumptions}
\newtheorem{lem}{Lemma}
%\crefname{lem}{Lemma}{Lemmas}
\newtheorem{thm}{Theorem}
\newtheorem{prop}{Proposition}
\newtheorem{cor}{Corollary}
\newtheorem{claim}{Claim}

\DeclareMathOperator*{\argmax}{arg\,max}
\DeclareMathOperator*{\argmin}{arg\,min}
\DeclareMathOperator*{\tr}{tr}
\DeclareMathOperator*{\rk}{rk}
\DeclareMathOperator*{\spa}{span}
\DeclareMathOperator*{\ve}{vec}
\DeclareMathOperator*{\diag}{diag}
\DeclareMathOperator*{\Rea}{Re}
\DeclareMathOperator*{\var}{Var}
\DeclareMathOperator*{\cov}{Cov}
\DeclareMathOperator*{\corr}{Corr}

\setlength{\skip\footins}{\baselineskip}
\setlength{\footnotesep}{0.3\baselineskip}

\allowdisplaybreaks

\begin{document}

{\noindent \bf University of Chicago \hfill
Spring 2026 \\
Department of Economics \hfill ECMA 31150}

\begin{center}
{\Large Problem Set 3} \\
This version: \today \\[1ex]
Mikkel Plagborg-Moller \\
\vspace{0.5\baselineskip}
Due Friday, May 1
\end{center}
\vspace{0.2\baselineskip}

\noindent \emph{Please upload your answers on Canvas, with code attached. You are encouraged to collaborate, but please write up your answers independently and do not share code chunks. Use any software you like, although you learn the most by limiting the use of generative AI. If you find any errors or require clarification, please let me know right away.}

\subsection*{Question 1}
Consider the following model for $t=1,2,\dots,T$:
\begin{align*}
y_t &= x_t + u_t + u_{t-1}, \\
x_t &= \tfrac{1}{2} x_{t-1} + w_t.
\end{align*}
The processes $\lbrace u_t \rbrace_{t=0}^\infty$ and $\lbrace w_t \rbrace_{t=1}^\infty$ are independent of each other, and each is i.i.d.\ standard normally distributed. Finally, $x_0 \sim N(0,4/3)$, independently of $\lbrace u_t,w_t \rbrace$.

Let
\[\hat{\beta} \equiv \frac{\sum_{t=1}^T y_tx_t}{\sum_{t=1}^T x_t^2}\]
be the OLS estimator in a regression of $y_t$ on $x_t$ (without intercept).

\begin{enumerate}[i)]
\item Show that $\hat{\beta} \overset{p}{\to} \beta$ as $T \to \infty$ for some constant $\beta$. What is $\beta$?
\item Derive the asymptotic distribution of $\sqrt{T}(\hat{\beta}-\beta)$. Calculate the explicit numerical values of the parameters of the asymptotic distribution.
\end{enumerate}
[Note: when applying relevant LLNs or CLTs, you do not need to verify the technical regularity conditions for these theorems.]

\subsection*{Question 2}
We seek to estimate the consumption Euler equation
\[y_t = c + \delta_f E_t[y_{t+1}] + \delta_b y_{t-1} - \tau(i_t - E_t[\pi_{t+1}]) + u_t.\]
Assume that $u_t$ is a demand shock that is unpredictable from past data. Assume also that expectations are rational. In the following we will use output (GDP) in place of consumption, as in the paper by Fuhrer \& Rudebusch (2004) on the reading list.

\begin{enumerate}[i)]
\item Define $r_t \equiv i_t-\pi_{t+1}$. Argue that the parameters of the Euler equation can be estimated via an IV regression of $y_t$ on the covariates $(y_{t+1}, y_{t-1}, r_t)$, using lagged data as instruments.
\end{enumerate}
Go to the St.\ Louis Fed's FRED website and download the series for real GDP (FRED code: GDPC1), the consumer price index (code: CPIAUCSL), and the federal funds rate (code: FEDFUNDS). Transform both monthly series into quarterly averages. Then transform the quarterly CPI series to annualized quarter-over-quarter log inflation: $\pi_t \equiv 400 \times [\log(\text{CPI}_t)-\log(\text{CPI}_{t-1})]$. Use the largest overlapping sample for the data, up to and including the end of 2019.

Theoretically, the variable $y_t$ should be measured in deviations from a long-run trend. Run a regression of $\log(\text{real GDP})$ on a linear and a quadratic function of time $t$ (and an intercept). Treat the residuals from this regression as the $y_t$ that enters into the Euler equation (multiply the residuals by 400 to make the units comparable to $i_t$ and $\pi_{t+1}$). Finally, we identify $i_t$ with the (quarterly average of the) federal funds rate.
\begin{enumerate}[i)]
\setcounter{enumi}{1}
\item Estimate the parameters of the Euler equation using an IV regression with HAC standard errors. Use two lags of $y_t$, $\pi_t$, and $i_t$ as instruments.
\item Is $\tau$ statistically significant, using a 10\% significance level?
\item Do the instruments appear to be weak?
\end{enumerate}

\clearpage

\subsection*{Question 3}
The course website contains a data file \texttt{oil\_shock.csv} with an updated version of the oil news shock series constructed in the paper by K\"{a}nzig (2021) on the reading list.\footnote{This series was downloaded from \url{https://github.com/dkaenzig/oilsupplynews}} Download additionally the following series from FRED: industrial production (FRED code: INDPRO), the federal funds rate (code: FEDFUNDS), the consumer price index (code: CPIAUCSL), and the West Texas Intermediate spot crude oil price (code: WTISPLC). Transform industrial production, CPI, and the WTI oil price to logs prior to the analysis. Use the largest available overlapping sample for the five data series (shock and four FRED series), up to and including the end of 2019.

\begin{enumerate}[i)]
\item Using local projections, estimate the impulse response function of (log) CPI with respect to the oil news shock at horizons $h=0,1,\dots,48$. Control for lags of all five variables mentioned above, with the number of lags selected by AIC. Plot the estimated IRF along with a 68\% confidence band.
\item Interpret verbally the estimated impulse response at horizon 24, i.e., discuss its sign, magnitude, and units from an economic point of view. Note that the oil news shock variable is measured in percent.
\item Read the introduction and Section I of K\"{a}nzig (2021). Discuss in your own words whether you believe the oil news shock is a well-identified ``shock''. Provide some statistical and/or visual diagnostics to support your discussion. This is an open-ended question, but please keep your answer concrete and to the point.
\end{enumerate}

\end{document}
